\chapter{Principal Component Analysis }
\label{ch:pca}


Machine learning is the field of study that gives computers the ability to learn from data without being explicitly programmed. In this context, a model is a mathematical function or computational structure that transforms input data into an output, following rules that are not explicitly defined but inferred from examples. The data are represented through numerical features, and the model uses them to learn relationships, trends or structures that are relevant for a given task, such as classification or dimensionality reduction. During training, the model adjusts its internal parameters according to the examples it receives, in order to capture patterns that remain informative when new observations are presented.

This approach is well suited to Raman spectroscopy of biological tissues. As discussed in \Cref{ch:introduction}, the complexity of biological Raman spectra makes it difficult to reduce the classification problem to a few manually selected bands. For this reason, the entire spectrum is treated as a numerical object whose overall profile may contain information relevant to the classification task. In the dataset used in this work, the Raman measurements are already organized in a computationally suitable tabular form. Starting from the structure described in \Cref{ch:Dataset}, the input data for the machine learning analysis are obtained from the spectral columns \texttt{band\_}n. These columns contain the Raman intensity values at the 483 selected shifts; from a machine learning perspective, they form the feature vector associated with each pixel-level spectrum acquired from a Raman map. When all spectra are considered together, the corresponding feature vectors are arranged into a data matrix, whose rows represent spectra and whose columns represent Raman-shift variables. 

Machine learning methods are commonly distinguished according to the information available during training. When the aim is to train a classifier, the learning process is supervised: the model is shown both the input data and the correct output associated with each example. In this way, it can learn how spectral information is related to tissue classes and use this relation to assign a class to new spectra. In the present work, this corresponds to using each Raman spectrum as input and its associated tissue label as the target output for binary classification.

When labels are not part of the learning process, the analysis becomes unsupervised: the algorithm only receives the input data and searches for structure within them. This approach is useful when the goal is not to assign a class directly, but to explore the organization of the data, visualize similarities between spectra or reduce the number of variables. PCA belongs to this second case. 

This chapter focuses on Principal Component Analysis (PCA), an unsupervised dimensionality-reduction method that identifies the directions along which the spectra show the largest variance. PCA projects the original spectral variables onto a new set of components, providing a more compact representation of the data. This transformation is useful because dimensionality reduction can be performed by retaining only the first principal components, which account for most of the spectral variance, and discarding later components associated with smaller contributions. In this way, each spectrum is represented by a reduced set of PCA scores instead of the full set of original Raman-shift variables. This compact representation can reduce redundancy and may limit the influence of weak or noisy variations before classification.

In this work, PCA is therefore first used as an
exploratory analysis of the dataset. Its theoretical formulation is presented with particular attention to the identification of the principal components, the amount of variance retained after projection, and the interpretation of the spectra in the transformed PCA space. When PCA is later used inside the classification pipeline, however, it must be fitted only on the training data of each split or fold and then applied to the corresponding validation or test data. This preserves the separation between training and evaluation data, preventing information from the validation or test set from influencing the learned representation.
